\section{Methodology}
\label{sec:methods}
\subsection{End-to-End Pipeline and Data Loading}
The pipeline uses a custom subset wrapper and collation function around torchvision's
raw Fashion-MNIST dataset. Training and validation wrappers use disjoint official-training
indices and independent transform objects; the test wrapper uses the official test
source. Training shuffles through a dedicated run-seeded generator. Validation and
test do not shuffle. All loaders retain IDs and use \texttt{drop\_last=False}.

Figure~\ref{fig:pipeline} separates the update and evaluation paths. Training enables
dropout and gradients; validation/test use evaluation mode and inference without
gradient tracking. Both consume the same image/label batch contract.
\begin{figure}[H]\centering
\begin{tikzpicture}[>=Latex,node distance=8mm and 9mm,
 box/.style={draw,rounded corners,align=center,text width=39mm,minimum height=12mm,
 font=\small,fill=blue!5}]
\node[box] (raw) {Raw IDX + fixed\\split manifest};
\node[box,right=of raw] (prep) {Train crop / eval\\tensor + normalization};
\node[box,right=of prep] (loader) {DataLoader\\images, labels, IDs};
\node[box,below=of prep] (model) {Model forward\\raw logits \([B,10]\)};
\node[box,below left=of model,fill=orange!10] (loss) {Cross-entropy\\backward + Adam};
\node[box,below right=of model,fill=green!8] (eval) {Argmax + probabilities\\metrics + saved outputs};
\draw[->] (raw)--(prep); \draw[->] (prep)--(loader);
\draw[->] (loader.south)|-(model.east);
\draw[->] (model.south west)--node[midway,above left,font=\small]{train} (loss.north east);
\draw[->] (model)--node[above right,font=\small]{evaluate} (eval);
\draw[->,dashed] (loss.west)-|([xshift=-5mm]raw.west)|-
node[pos=.8,above,font=\small]{parameter update} (model.west);
\end{tikzpicture}
\caption{Data preparation and model interface, with distinct update/evaluation paths.
Softmax probabilities support analysis; cross-entropy receives raw logits.}
\label{fig:pipeline}
\end{figure}
The engine computes sample-weighted loss and epoch-level accuracy/F1 from online
training predictions. These include augmentation and training-mode behavior.
Thus, close train/validation curves alone do not prove absence of overfitting.

\subsection{Shared Model Interface and Input Representations}
Every classifier accepts float32 \([B,1,28,28]\) and returns raw \([B,10]\) logits.
Representation changes occur inside the model, leaving data loading and metrics
unchanged. Table~\ref{tab:representations} connects each representation to its
structural preference.
\begin{table}[H]\centering\small
\caption{Representations used in the reported configurations.}
\label{tab:representations}
\begin{tabularx}{\linewidth}{lYY}
\toprule Model & Internal representation & Structural preference\\\midrule
Linear & Flattened 784-vector & One affine decision rule per class\\
MLP & Flattened vector, dense hidden features & Nonlinear global combinations\\
CNN & Two-dimensional feature maps & Local connectivity and shared filters\\
GRU & 28 row timesteps of width 28 & Ordered recurrent context\\
Transformer & 49 \(4\times4\) patches, width 64 & Position-aware token interactions\\
\bottomrule
\end{tabularx}
\end{table}
The input shape and label space are shared despite different learned representations
and parameter counts.

\subsection{Model Architectures}
\subsubsection{Linear / Softmax Classifier}
The baseline flattens each image and applies one affine layer:
\[
 x=\operatorname{flatten}(X)\in\mathbb{R}^{784},\qquad
 z=Wx+b,\quad W\in\mathbb{R}^{10\times784}.
\]
There is no hidden nonlinearity. The \(784\times10+10=7,850\) parameters are all
trainable; flattening adds none. Probabilities, when needed, are
\(p_k=\exp(z_k)/\sum_{j=0}^{9}\exp(z_j)\).
The implementation returns logits because cross-entropy already combines
log-softmax and the label loss.\citep{pytorchce}
The project implements the wrapper and validation around PyTorch
\texttt{nn.Flatten} and \texttt{nn.Linear}. This small baseline has neither
a hidden feature hierarchy nor built-in image locality.

\subsubsection{Multilayer Perceptron}
The MLP uses \(784\rightarrow256\rightarrow128\rightarrow10\).
Every hidden affine layer is followed by ReLU and dropout 0.1; the final layer
returns logits. For hidden layer \(\ell\),
\[
 a_\ell=W_\ell h_{\ell-1}+b_\ell,\qquad
 h_\ell=\operatorname{ReLU}(a_\ell).
\]
ReLU keeps positive values and sets negative values to zero, enabling nonlinear
boundaries. Without nonlinearities, consecutive affine layers collapse into one
affine transformation. Training-time dropout masks hidden activations and scales
retained values; evaluation disables these masks. It regularizes dense features,
but does not favor spatially adjacent pixels.
The parameter count is
\[
 (784\cdot256+256)+(256\cdot128+128)+(128\cdot10+10)=235,146.
\]
Project code assembles the network from PyTorch layers. Other activation names
are supported, but only ReLU is used in the reported MLP runs; no activation
ablation is claimed.

\subsubsection{Custom Convolutional Neural Network}
Convolutions apply learned filters to local neighborhoods, reusing weights across
positions. Multiple filters yield feature maps responsive to patterns learned
from data. For stride one, a filter response has the form
\[
 Z_{m,i,j}=b_m+\sum_c\sum_u\sum_v W_{m,c,u,v}X_{c,i+u,j+v}.
\]
For illustration, a \(2\times2\) filter with ones on its main diagonal gives
\(1+6=7\) on \(\bigl[\begin{smallmatrix}1&2\\5&6\end{smallmatrix}\bigr]\).
This is a numerical example; the actual network uses \(3\times3\) kernels.

Max pooling retains the largest value in a region: the same example pools to 6,
versus 3.5 under average pooling. Pooling reduces spatial resolution and may reduce
sensitivity to small shifts, while discarding fine spatial detail. Its aggregation
rule has no trainable weights.

Table~\ref{tab:cnn} gives the actual custom architecture. Padding preserves each
convolution's size; max pooling reduces spatial dimensions \(28\to14\to7\).
The implementation uses \texttt{ceil\_mode=True}; these even-sized inputs have
the dimensions shown. Adaptive average pooling gives one value per channel.
\begin{table}[H]\centering\small
\caption{Custom CNN layers and resulting tensor shapes.}
\label{tab:cnn}
\begin{tabular}{ll}
\toprule Layer & Output\\\midrule
Input & \([B,1,28,28]\)\\
Conv \(1\to32\), \(3\times3\), pad 1; ReLU & \([B,32,28,28]\)\\
MaxPool \(2\times2\) & \([B,32,14,14]\)\\
Conv \(32\to64\), \(3\times3\), pad 1; ReLU & \([B,64,14,14]\)\\
MaxPool \(2\times2\) & \([B,64,7,7]\)\\
AdaptiveAvgPool; flatten & \([B,64]\)\\
Linear \(64\to128\); ReLU; Dropout 0.1 & \([B,128]\)\\
Linear \(128\to10\) & \([B,10]\)\\\bottomrule
\end{tabular}
\end{table}
All 28,426 parameters are trainable. The project defines this network using
PyTorch convolution, activation, pooling and linear layers; it is not a pretrained
CNN backbone.

\subsubsection{GRU-Based Image Classifier}
Each image row is one timestep,
\[
 x_t=X_{t,:}\in\mathbb{R}^{28},\qquad t=1,\ldots,28.
\]
The batch-first sequence has shape \([B,28,28]\). A single unidirectional
\texttt{nn.GRU} uses input width 28 and hidden width 128; its final hidden
state is projected to ten logits. A timestep is a complete row, not one pixel.
The traversal carries information from earlier rows into later decisions.

The equations below match the PyTorch cell actually used:\citep{pytorchgru}
\begin{align*}
 r_t&=\sigma(W_{ir}x_t+b_{ir}+W_{hr}h_{t-1}+b_{hr}),\\
 z_t&=\sigma(W_{iz}x_t+b_{iz}+W_{hz}h_{t-1}+b_{hz}),\\
 n_t&=\tanh(W_{in}x_t+b_{in}+r_t\odot(W_{hn}h_{t-1}+b_{hn})),\\
 h_t&=(1-z_t)\odot n_t+z_t\odot h_{t-1}.
\end{align*}
The reset gate modulates transformed past information in the candidate; the update
gate blends it with retained memory. PyTorch multiplies by the reset gate after
the hidden affine transformation.

Recurrent weights and two bias vectors contribute
\(3(28\cdot128+128\cdot128+2\cdot128)=60,672\) parameters.
The head adds 1,290, totaling 61,962. The project wrapper handles row conversion
and state selection, while the cell is a PyTorch component. Although the code
supports columns, LSTM and bidirectionality, these variants were not evaluated.
The reported single-layer GRU has no inter-layer dropout transition.

\subsubsection{Custom Transformer Classifier}
The custom Transformer divides an image into a \(7\times7\) grid of
non-overlapping \(4\times4\) patches. The 49 patches are flattened to 16 values
each and projected to embedding width 64. A learned \(49\times64\) position
embedding is added:
\[
 P\in\mathbb{R}^{B\times49\times16},\qquad
 H_0=PW_E+b_E+E_{\rm pos}\in\mathbb{R}^{B\times49\times64}.
\]
This represents local regions as tokens while preserving their image positions.
Two pre-normalized residual blocks contain four-head self-attention and an FFN.
Each head has \(d_h=16\):
\[
 Q=HW_Q,\quad K=HW_K,\quad V=HW_V,\qquad
 A=\operatorname{softmax}\!\left(QK^\mathsf{T}/\sqrt{16}\right)V.
\]
Per image, attention weights have shape \([4,49,49]\). Each token attends to all
49 tokens and receives a 16-value result per head; concatenation and an output
projection restore width 64. The mechanism follows scaled dot-product
attention,\citep{attention} with the project's residual structure
\begin{align*}
 U_\ell&=H_{\ell-1}+
 \operatorname{Dropout}\bigl(\operatorname{MHSA}(\operatorname{LN}(H_{\ell-1}))\bigr),\\
 H_\ell&=U_\ell+\operatorname{FFN}(\operatorname{LN}(U_\ell)).
\end{align*}
The FFN is \(64\to256\to64\), with GELU and dropout 0.1. After two blocks,
final layer normalization and mean pooling over 49 tokens give a 64-value vector;
a \(64\to10\) head returns logits. There is no class token.

Parameter counts are 1,088 for patch projection, 3,136 for positions,
two blocks of 49,984, 128 for final normalization and 650 for the head:
104,970 in total. The project implements patchification, Q/K/V projections,
attention, residual blocks and pooling with PyTorch tensor operations and layers.
It does not substitute a built-in Transformer encoder or pretrained vision model.
\FloatBarrier

\subsection{Loss Function and Evaluation Metrics}
Every model uses unweighted multiclass cross-entropy:
\[
 L=-\frac{1}{N}\sum_{i=1}^{N}
 \log\frac{\exp(z_{i,y_i})}{\sum_{k=0}^{9}\exp(z_{i,k})}.
\]
Loss is aggregated by sample count, correctly weighting the incomplete final
batch. Argmax logits give labels. For a confusion matrix whose rows are true
classes and columns predictions,
\[
 \mathrm{Accuracy}=\frac{\sum_c C_{cc}}{\sum_{c,k}C_{ck}},\qquad
 P_c=\frac{TP_c}{TP_c+FP_c},\quad R_c=\frac{TP_c}{TP_c+FN_c},
\]
\[
 F1_c=\frac{2P_cR_c}{P_c+R_c},\qquad
 \mathrm{MacroF1}=\frac{1}{10}\sum_{c=0}^{9}F1_c.
\]
Undefined class precision, recall or F1 is assigned zero, matching
\texttt{zero\_division=0}. Macro-F1 weights classes equally and can reveal
weaknesses hidden by overall accuracy. Tables use the mean and sample SD across
three training seeds, with denominator \(n-1\). Scores converted to percent are
labeled explicitly; differences between percent scores are percentage points.

\subsection{Training, Hyperparameter Selection and Checkpointing}
\label{sec:training}
All 15 main runs use Adam, learning rate \(10^{-3}\), weight decay 0, batch size
64 and a 100-epoch cap. The bundle contains one effective configuration per model
repeated across three seeds; it does not document an exhaustive hyperparameter
search. Momentum fields present in Linear/MLP configs do not change Adam into SGD.

Linear and MLP use cosine annealing over 100 epochs to a minimum learning rate
\(10^{-5}\), with early-stopping patience 15 and macro-F1 minimum improvement
0.001. CNN, GRU and Transformer have no scheduler and early stopping disabled.
Equal epoch caps therefore do not imply equal completed budgets or training policies.

The selected checkpoint maximizes validation macro-F1, breaking exact ties by
lower validation loss and earlier epoch. Test results never enter this choice.
The early-stopping counter is separate: for Linear/MLP it resets only after an
improvement larger than 0.001. A smaller improvement may update \texttt{best.pt}
without resetting patience.

The trainer saves configs, source snapshots, metadata, environment records,
history, selected metrics and checkpoints. Evaluation uses \texttt{best.pt}.
The last checkpoint retains optimizer, scheduler and RNG state for supported
epoch-boundary resumption; no mid-minibatch resume guarantee is claimed.

\subsection{Experimental Design and Fairness Controls}
\begin{table}[H]\centering\small
\caption{Questions, controls and interpretation limits.}
\label{tab:design}
\begin{tabularx}{\linewidth}{>{\raggedright\arraybackslash}p{25mm}YYY}
\toprule Focus & Changed factors & Fixed controls & Evidence/criterion\\\midrule
Quality & Model configuration and representation & Split, preprocessing, labels, evaluator & Accuracy, macro-F1 and CE across three seeds\\
Representation & Flat pixels, spatial maps, rows, patches & Image task and held-out test & Class errors; association, not isolated causality\\
Efficiency & Architecture, completed epochs, kernels & Batch 64 within each timing protocol & Parameters, GPU loop time, separate CPU forward time\\
Seed variation & Run seed & Per-model effective config & Mean/sample SD and paired seed differences\\
\bottomrule
\end{tabularx}
\end{table}
This is a shared-data comparison, not matched-parameter or matched-compute
experimentation. No row/column, LSTM/GRU, dropout or augmentation ablation is
claimed. Such follow-ups could test proposed explanations; their outcomes
are not established by this comparison.

\subsection{Environment, Timing and Reproduction Protocol}
\label{sec:environment}
Training used Google Colab Linux and an NVIDIA Tesla T4 (14.56 GiB recorded memory),
PyTorch 2.11.0+cu130, torchvision 0.26.0+cu130, NumPy 2.1.3 and CUDA 13.0.
Python is 3.13.15 for 13 runs and 3.13.16 for Transformer seeds 67/69.
Inputs and models use float32 without AMP. The saved settings are matmul precision
\texttt{highest}, CUDA matmul TF32 disabled, cuDNN TF32 allowed and cuDNN
benchmarking disabled; deterministic algorithms are enabled. Python patch-version
variation is a reproducibility limitation, not a verified cause of score differences.

Validation replay and final-test evaluation used a separate Colab CPU session:
Python 3.13.16, PyTorch 2.11.0+cpu, torchvision 0.26.0+cpu, NumPy 2.1.3 and two
threads. Replay accuracy/F1 match the original validation scores in all runs;
the largest loss difference is \(3.84\times10^{-8}\).

Reported inference speed uses the separately retained \textbf{Windows CPU profile},
measured on 24 October 2026: Python 3.14.6, PyTorch 2.13.0+cpu, torchvision
0.28.0+cpu, NumPy 2.5.2 and two threads. The recorded processor is Intel64
Family 6 Model 170 Stepping 4, GenuineIntel (AMD64); no friendly product name
was recorded. These are not Colab or T4 speed measurements.

For each selected checkpoint, the forward-only protocol uses a preallocated
zero tensor \([64,1,28,28]\), float32, evaluation mode, 20 warm-ups and 100 timed
forwards. It excludes checkpoint I/O, loading and transfer. Per-run latency is
the median batch time; amortized latency divides that median by 64; throughput
is 64 divided by the median in seconds. Tables average these run-level values
and give sample SD across three seeds. This is not batch-one or end-to-end speed.
\FloatBarrier
